{\itshape
(ii) Suppose that one has one of the following two properties:
\begin{enumeratea}
\item \(X\) is an \(S\)-group smooth over \(S\);
\item \(X=\Aut_{S}(E)\) where \(E\) is smooth and affine over \(S\).
\end{enumeratea}
\noindent Then \(f_{0}\) extends to a morphism of \(S\)-groups \(Y\to X\); two such
extensions are conjugate by an inner automorphism of \(X\) defined by a section
of \(X\) over \(S\) inducing the unit section of \(X_{0}\) over \(S_{0}\).
\par}

\begin{proof}
\renewcommand{\qedsymbol}{}
Let us introduce the \(S_{n}\) as above.%
{\renewcommand{\thefootnote}{(59)}\nde{In what follows, one has slightly
modified the original and set it out in detail.}}
For~(ii), let us first note that a scheme smooth over \(S\) is necessarily
locally of finite presentation over \(S\); hence, in case~(b), \(E\), being
smooth and affine over \(S\), is necessarily of finite presentation over \(S\),
\ie one is indeed under hypothesis~(b) of~2.0.

Then, under the hypotheses of~(ii), condition~\((i_{1})\) of~2.1 is
automatically satisfied by~0.16 and~0.17; moreover every section of
\(X_{S_{n}}\) over \(S_{n}\) lifts to a section of \(X_{S_{n+1}}\) over
\(S_{n+1}\), by the definition of ``smooth over \(S\)'' in case~(a), and
by~0.17 in case~(b). Consequently, if \(f\) and \(f'\) are two liftings of
\(f_{0}\), one may suppose step by step that \(f_{n}=f'_{n}\), by lifting the
inner automorphism whose existence is asserted by part~(ii) of the theorem,
which completes the proof.
\end{proof}

Arguing in the same way, one obtains, taking Remark~2.3 into account:

\begin{corollary}{2.6}\label{III.2.6}
Let \(S\) be a scheme, \(\mathcal{I}\) a nilpotent ideal defining the closed
subscheme \(S_{0}\), \(Y\) an \(S\)-group flat over \(S\) and affine, \(X\) an
\(S\)-group smooth over \(S\).
\begin{enumeratei}
\item If, for every \(n\geqslant 0\), one has
\(H^{2}\bigl(Y_{0},\Lie(X_{0}/S_{0})\otimes_{\mathcal{O}_{S_{0}}}
\mathcal{I}^{n+1}/\mathcal{I}^{n+2}\bigr)=0\), every morphism of \(S_{0}\)-groups
\(f_{0}\colon Y_{0}\to X_{0}\) lifts to a morphism of \(S\)-groups
\(f\colon Y\to X\).
\item If, for every \(n\geqslant 0\), one has
\(H^{1}\bigl(Y_{0},\Lie(X_{0}/S_{0})\otimes_{\mathcal{O}_{S_{0}}}
\mathcal{I}^{n+1}/\mathcal{I}^{n+2}\bigr)=0\), two such liftings are conjugate
by an inner automorphism of \(X\) defined by a section of \(X\) over \(S\)
inducing the unit section of \(X_{0}\) over \(S_{0}\).
\end{enumeratei}
\end{corollary}

Now one has the following lemma:

\begin{lemma}{2.7}\label{III.2.7}
Let \(S\) be an affine scheme, \(G\) an affine \(S\)-group, \(\mathcal{F}\) a
quasi-coherent \(\mathcal{O}_{S}\)-module, \(\mathcal{L}\) a locally free
\(\mathcal{O}_{S}\)-module. Suppose that one has an operation of \(G\) on
\(\mathcal{F}\) in the sense of Exposé~I, which defines an operation of \(G\)
on \(\mathcal{F}\otimes_{\mathcal{O}_{S}}\mathcal{L}\)%
{\renewcommand{\thefootnote}{(60)}\nde{Where \(\mathcal{L}\) is equipped with
the trivial action of \(G\).}}.
Denote by \(\Lambda\) the ring of \(S\), and by \(L\) the \(\Lambda\)-module
defining \(\mathcal{L}\) (which is therefore a projective module). One has a
canonical isomorphism
\[
H^{*}(G,\mathcal{F}\otimes_{\mathcal{O}_{S}}\mathcal{L})
  \simeq
  H^{*}(G,\mathcal{F})\otimes_{\Lambda}L.
\]
\end{lemma}

\begin{proof}
\renewcommand{\qedsymbol}{}
{\renewcommand{\thefootnote}{(61)}\nde{One has set out in detail (and
simplified) the proof of the original (the latter moreover invoked the
isomorphisms
\(\mathcal{H}^{n}(\mathcal{C}\otimes_{\mathcal{O}_{S}}\mathcal{L})
\simeq\mathcal{H}^{n}(\mathcal{C})\otimes_{\mathcal{O}_{S}}\mathcal{L}\)
and \(H^{n}(\Gamma(-))\simeq\Gamma(\mathcal{H}^{n}(-))\), where the
\(\mathcal{H}^{n}(\mathcal{C})\) denote the cohomology sheaves of the complex
\(\mathcal{C}\)).}}
Indeed, denote by \(\mathcal{A}\) the \(\mathcal{O}_{S}\)-algebra
\(\mathcal{A}(G)\) and consider the complex \(\mathcal{C}\) of quasi-coherent
\(\mathcal{O}_{S}\)-modules:
\[
0\longrightarrow\mathcal{F}
  \longrightarrow\mathcal{F}\otimes_{\mathcal{O}_{S}}\mathcal{A}
  \longrightarrow\mathcal{F}\otimes_{\mathcal{O}_{S}}\mathcal{A}
    \otimes_{\mathcal{O}_{S}}\mathcal{A}
  \longrightarrow\cdots
\]
By~I, 5.3, \(H^{*}(G,\mathcal{F})\)
(\resp \(H^{*}(G,\mathcal{F}\otimes_{\mathcal{O}_{S}}\mathcal{L})\))
is the cohomology of the complex \(\Gamma(S,\mathcal{C})\)
(\resp \(\Gamma(S,\mathcal{C}\otimes_{\mathcal{O}_{S}}\mathcal{L})\)).
Now, as \(S\) is affine, one has (\cf EGA~I, 1.3.12)
\[
\Gamma(S,\mathcal{C}\otimes_{\mathcal{O}_{S}}\mathcal{L})
  \simeq
  \Gamma(S,\mathcal{C})\otimes_{\Lambda}L.
\]
As \(L\) is a projective \(\Lambda\)-module (hence flat), one also has
\(H^{*}\bigl(\Gamma(S,\mathcal{C})\otimes_{\Lambda}L\bigr)
\simeq H^{*}\bigl(\Gamma(S,\mathcal{C})\bigr)\otimes_{\Lambda}L\),
whence the announced result.
\end{proof}

Using the lemma, one transforms~2.6 into:

\begin{corollary}{2.8}\label{III.2.8}
Let \(S\) be an affine scheme, \(\mathcal{I}\) a nilpotent ideal on \(S\)
defining the closed subscheme \(S_{0}\). Suppose that the
\(\mathcal{I}^{n+1}/\mathcal{I}^{n+2}\) are locally free over \(S_{0}\). Let \(Y\)
be an \(S\)-group flat over \(S\) and affine, \(X\) an \(S\)-group smooth over
\(S\), and \(f_{0}\colon Y_{0}\to X_{0}\) a morphism of \(S\)-groups.
\begin{enumeratei}
\item If \(H^{2}\bigl(Y_{0},\Lie(X_{0}/S_{0})\bigr)=0\), \(f_{0}\) lifts to a
morphism of \(S\)-groups \(Y\to X\).
\item If \(H^{1}\bigl(Y_{0},\Lie(X_{0}/S_{0})\bigr)=0\), two such liftings are
conjugate by an inner automorphism of \(X\) defined by a section of \(X\) over
\(S\) inducing the unit section of \(X_{0}\) over \(S_{0}\).
\end{enumeratei}
\end{corollary}

In particular, setting \(Y=X\):

\begin{corollary}{2.9}\label{III.2.9}
Let \(S\) and \(S_{0}\) be as above. Let \(X\) be an \(S\)-group smooth over
\(S\) and affine.
\begin{enumeratei}
\item If \(H^{1}\bigl(X_{0},\Lie(X_{0}/S_{0})\bigr)=0\), every endomorphism of
\(X\) over \(S\) inducing the identity on \(X_{0}\) is the inner automorphism
defined by a section of \(X\) over \(S\) inducing the unit section of \(X_{0}\)
over \(S_{0}\).
\item If
\[
H^{2}\bigl(X_{0},\Lie(X_{0}/S_{0})\bigr)=0,
\]
every
\(S_{0}\)-automorphism of \(X_{0}\) extends to an \(S\)-automorphism of
\(X\).%
{\renewcommand{\thefootnote}{(62)}\nde{Indeed, let \(f_{0}\) be an
\(S_{0}\)-automorphism of \(X_{0}\) and let \(g_{0}\) be its inverse. By~2.8
(i), \(f_{0}\) (\resp \(g_{0}\)) lifts to an \(S\)-endomorphism
\(f\) (\resp \(g\)) of \(X\). Then \(g\circ f\) and \(f\circ g\)
are endomorphisms of \(X\) inducing the identity on \(X_{0}\); they are
therefore, by~0.7, \(S\)-automorphisms of \(X\), and the same is therefore true
of \(f\) and \(g\).}}
\end{enumeratei}
\end{corollary}

\begin{remark}{2.10}\label{III.2.10}
The assertions concerning the \(H^{1}\) have converses by the theorem. Let us
point out the following as an example: if \(S=\mathrm{I}_{S_{0}}\) is the
scheme of dual numbers over \(S_{0}\) (II, 2.1) and if \(X\) is a flat
\(S\)-group such that every automorphism of \(X\) over \(S\) inducing the
identity on \(S_{0}\) is the inner automorphism defined by a section of \(X\)
over \(S\) inducing the unit section of \(X_{0}\) over \(S_{0}\), then
\(H^{1}\bigl(X_{0},\Lie(X_{0}/S_{0})\bigr)=0\).%
% typo?: printed ``identité sur \(S_0\)''; kept as printed
{\renewcommand{\thefootnote}{(63)}\nde{This is used in XXIV, 1.13.}}
\end{remark}

\begin{corollary}{2.11}\label{III.2.11}
Let \(S\), \(\mathcal{I}\) and \(\mathcal{J}\) be as in~2.1. Let \(Y\) be an
\(S\)-group scheme flat over \(S\), \(X\) an \(S\)-group scheme,
\(f\colon Y\to X\) a morphism of \(S\)-groups. The set of morphisms from \(Y\)
to \(X\) deduced from \(f\) by conjugation by elements \(x\in X(S)\) inducing
the unit of \(X(S_{\mathcal{J}})\) is isomorphic to the quotient
\[
E=\Hom_{\mathcal{O}_{S_{0}}}\bigl(\omega^{1}_{X_{0}/S_{0}},\mathcal{J}\bigr)
  \big/
  {\Hom_{\mathcal{O}_{S_{0}}}\bigl(\omega^{1}_{X_{0}/S_{0}},\mathcal{J}\bigr)}
    ^{\mathrm{ad}(Y_{0})},
\]
where the second group is formed of the \(\mathcal{O}_{S_{0}}\)-morphisms
\(\omega^{1}_{X_{0}/S_{0}}\to\mathcal{J}\) which, by every base change
\(S'\to S_{0}\), give morphisms
\(\omega^{1}_{X_{0}/S_{0}}\otimes_{\mathcal{O}_{S_{0}}}\mathcal{O}_{S'}
\to\mathcal{J}\otimes_{\mathcal{O}_{S_{0}}}\mathcal{O}_{S'}\)
invariant under the action of \(Y_{0}(S')\) on the first factor.
\end{corollary}

\begin{proof}
\renewcommand{\qedsymbol}{}
By~2.1 (iii), one knows that the set sought is isomorphic to
\(\Gamma(L_{0})/H^{0}(Y_{0},L_{0})\). Now
\(\Gamma(L_{0})=\Hom_{\mathcal{O}_{S_{0}}}\bigl(\omega^{1}_{X_{0}/S_{0}},
\mathcal{J}\bigr)\) and \(H^{0}(Y_{0},L_{0})\) is evidently none other than
\(\Gamma(L_{0})^{\mathrm{ad}(Y_{0})}\) in the sense of the preceding statement.
\end{proof}

\begin{corollary}{2.12}\label{III.2.12}
Under the conditions of~2.11, suppose moreover \(\omega^{1}_{X_{0}/S_{0}}\)
locally free of finite rank. Then
\[
E\simeq\Gamma\bigl(S_{0},\Lie(X_{0}/S_{0})\otimes_{\mathcal{O}_{S_{0}}}
\mathcal{J}\bigr)
  \big/
  H^{0}\bigl(Y_{0},\Lie(X_{0}/S_{0})\otimes_{\mathcal{O}_{S_{0}}}\mathcal{J}\bigr).
\]
\end{corollary}

\begin{proof}
\renewcommand{\qedsymbol}{}
{\renewcommand{\thefootnote}{(64)}\nde{One has added the sentence which
follows.}}
Indeed, if \(\omega^{1}_{X_{0}/S_{0}}\) is locally free of finite rank, one has
\(\mathcal{H}\mathrm{om}_{\mathcal{O}_{S_{0}}}\bigl(\omega^{1}_{X_{0}/S_{0}},
\mathcal{J}\bigr)
\simeq\Lie(X_{0}/S_{0})\otimes_{\mathcal{O}_{S_{0}}}\mathcal{J}\).
\end{proof}

\begin{corollary}{2.13}\label{III.2.13}
Suppose moreover \(Y_{0}\) diagonalizable. Then
\[
E\simeq\Gamma\bigl(S_{0},\Lie(X_{0}/S_{0})\otimes_{\mathcal{O}_{S_{0}}}
\mathcal{J}\bigr)
  \big/
  \Gamma\bigl(S_{0},\Lie(X_{0}/S_{0})^{\mathrm{ad}(Y_{0})}
    \otimes_{\mathcal{O}_{S_{0}}}\mathcal{J}\bigr)
\]
where \(\Lie(X_{0}/S_{0})^{\mathrm{ad}(Y_{0})}\) may be constructed as the
factor of the decomposition of~I, 4.7.3, corresponding to the zero character
of \(Y_{0}\).
\end{corollary}

\begin{proof}
\renewcommand{\qedsymbol}{}
Indeed, if \(Y_{0}\simeq\mathrm{D}_{S_{0}}(M)\), one has by \loccit a
decomposition into a direct sum:
\[
\Lie(X_{0}/S_{0})
  =\Lie(X_{0}/S_{0})_{0}
  \oplus
  \bigoplus_{\substack{m\in M\\ m\neq 0}}\Lie(X_{0}/S_{0})_{m}.
\]
On tensoring by \(\mathcal{J}\), one finds an analogous decomposition for
\(\Lie(X_{0}/S_{0})\otimes_{\mathcal{O}_{S_{0}}}\mathcal{J}\), whence the
relation
\[
H^{0}\bigl(Y_{0},\Lie(X_{0}/S_{0})\otimes\mathcal{J}\bigr)
  \simeq
  \Gamma\bigl(S_{0},\Lie(X_{0}/S_{0})_{0}\otimes_{\mathcal{O}_{S_{0}}}
    \mathcal{J}\bigr).
\]
\end{proof}

\begin{corollary}{2.14}\label{III.2.14}
Suppose moreover \(S_{0}\) affine. Then
\[
E\simeq\Gamma\Bigl(S_{0},
  \bigl[\Lie(X_{0}/S_{0})\big/\Lie(X_{0}/S_{0})^{\mathrm{ad}(Y_{0})}\bigr]
  \otimes_{\mathcal{O}_{S_{0}}}\mathcal{J}\Bigr).
\]
\end{corollary}

\sgasection{3}{Infinitesimal extensions of a group scheme}
\label{III.sec.3}

Still in the notations of \No 0 (\(S\), \(\mathcal{I}\), \(\mathcal{J}\),
etc.), let us give ourselves an \(S\)-scheme \(X\) and suppose \(X_{\mathcal{J}}\)
endowed with a group structure. We propose to find the \(S\)-group structures
on \(X\) inducing on \(X_{\mathcal{J}}\) the given structure.

From now on, we suppose \(X\) flat over \(S\). Let \(\mathcal{C}\) be the
category of \(S\)-schemes flat over \(S\). One therefore has
\(X\in\Ob\mathcal{C}\). We shall denote by \(Y\), \resp \(M\), the functor on
\(\mathcal{C}\) defined by \(X^{+}\), \resp \(L'_{X}\). The canonical morphism
\(p_{X}\colon X\to X^{+}\) defines a morphism of \(\widehat{\mathcal{C}}\)
\[
p\colon X\longrightarrow Y
\]
and the operation of \(L'_{X}\) on \(X\) in \(\widehat{\Sch}_{/S}\) defines an
operation of \(M\) on \(X\) in \(\widehat{\mathcal{C}}\). One checks at once
that \(X\) does indeed thus become formally principal homogeneous under
\(M_{Y}\) over \(Y\) (\cf 0.2 (i) and~0.4).

The operation of \(X^{+}\) on \(L'_{X}\) defined in~0.8 (denoted \(\Ad\) in
\loccit) defines an operation, denoted \(f\), of \(Y\) on \(M\). One knows, on
the other hand (0.5), that
\[
\Hom_{\widehat{\mathcal{C}}}(Z,M)\simeq\Hom_{S_{0}}(Z_{0},L_{0}),
  \qquad Z\in\Ob\mathcal{C},
\]
where \(L_{0}\) is the functor defined in~0.5.

\begin{lemma}{3.1}\label{III.3.1}
\begin{enumeratei}
\item The condition \((+)_{n}\) of~1.3 is satisfied for every positive integer
\(n\).
\item If one lets the \(S_{0}\)-group \(X_{0}\) act on the \(S_{0}\)-functor
\(L_{0}\) via its adjoint representation, one has a canonical isomorphism
\[
H^{*}(X_{0},L_{0})\simeq H^{*}(Y,M),
\]
(the first cohomology is computed in \(\Sch_{/S_{0}}\), the second in
\(\mathcal{C}\)).
\end{enumeratei}
\end{lemma}

\begin{proof}
\renewcommand{\qedsymbol}{}
The two parts of the lemma result from the relation:
\begin{align*}
\Hom_{\widehat{\mathcal{C}}}(Y,M)
  &\simeq\Hom_{\widehat{\Sch}_{/S_{0}}}\bigl(X^{+}\times_{S}S_{0},L_{0}\bigr) \\
  &\simeq\Hom_{S_{0}}(X_{0},L_{0}) \\
  &\simeq\Hom_{\widehat{\mathcal{C}}}(X,M),
\end{align*}
which comes at once from the definition of \(M\) as a
\(\prod_{S_{0}/S}\). This relation being more generally satisfied on replacing
\(X\), \(Y\) by \(X^{n}\), \(Y^{n}\), one deduces from it that every morphism
\(X^{n}\to M\) factors uniquely through \(Y^{n}\), which entails \((+)_{n}\).
One also deduces from it the relation \(C^{*}(Y,M)=C^{*}(X_{0},L_{0})\), which
entails~(ii).
\end{proof}

We may therefore apply the constructions of~1.3. In particular:

\begin{lemma}{3.2}\label{III.3.2}
Let \(P\colon X\times_{S}X\to X\) be a morphism. For \(P\) to induce the group
law of \(X_{\mathcal{J}}\), it is necessary and sufficient that \(P\) be an
admissible composition law (\cf 1.3.2) on \(X\).
\end{lemma}

\begin{proof}
\renewcommand{\qedsymbol}{}
Indeed, for \(P\) to induce the group law of \(X_{\mathcal{J}}\), it is
necessary and sufficient that \(P\) lift the group law of \(X^{+}\), or again
that of \(Y\). One therefore only has to show that every morphism \(P\) lifting
the group law of \(X_{\mathcal{J}}\) satisfies the identity \((++)\) of~1.3.2
(ii), which is exactly what one saw in~0.8.
\end{proof}

\begin{proposition}{3.3}\label{III.3.3}
Let \(S\) be a scheme and \(S_{0}\) a closed subscheme defined by a nilpotent
ideal. Let \(X\) be an \(S\)-scheme flat, and quasi-compact or locally of finite
presentation over \(S\). Let \(P\colon X\times_{S}X\to X\) be a composition law
on \(X\). For \(P\) to be a group law, it is necessary and sufficient that the
following two conditions be satisfied:
\begin{enumeratei}
\item \(P\) is associative.
% typo?: printed \(X_0=X\times_X S_0\); kept as printed
\item \(P\) induces on \(X_{0}=X\times_{X}S_{0}\) a group law.
\end{enumeratei}
\end{proposition}

\begin{proof}
\renewcommand{\qedsymbol}{}
These conditions are evidently necessary. Let us show that they are sufficient.
Suppose first that \(X\to S\) has a section. As \(X(S')\) is then nonempty
for each \(S'\to S\), it suffices%
{\renewcommand{\thefootnote}{(65)}\nde{One has corrected the original, by
removing the inadequate reference to an exercise of Bourbaki on semigroups
(\cf [BAlg], \S~I.2, Exercises~9 to~13) and by indicating the role of the
translations on the left and on the right; see the following N.D.E.}}
to show that, for every \(x\in X(S')\), the translations \emph{on the left and
on the right} by \(x\) are isomorphisms of \(X_{S'}\).%
{\renewcommand{\thefootnote}{(66)}\nde{Let \(E\) be a nonempty set equipped
with an associative composition law, such that every left translation
\(\ell_{x}\) is bijective; let us fix \(x_{0}\in E\). There exists a unique
\(e\in E\) such that \(x_{0}\cdot e=x_{0}\); then
\(x_{0}\cdot e\cdot x=x_{0}\cdot x\) entails \(e\cdot x=x\), for every
\(x\in E\). On the other hand, for every \(x\) there exists a unique \(x'\)
such that \(x\cdot x'=e\). Suppose moreover that there exists \(b\in E\) such
that the right translation \(r_{b}\) is injective. Then, for every \(x\), the
equality \(x\cdot e\cdot b=x\cdot b\) gives \(x\cdot e=x\) (\ie \(e\) is a
neutral element), and \(x\cdot x'\cdot x=x=x\cdot e\) entails \(x'\cdot x=e\),
\ie \(x'\) is the inverse of \(x\) on the left and on the right, hence \(E\) is
a group.\endgraf
Note that the hypothesis ``\(r_{b}\) injective'' is necessary: on every set
\(E\) one may define a composition law by \(x\cdot y=y\), for all
\(x,y\in E\); then every left translation is the identity (whence the
associativity of the law), but for every \(y\) one has \(r_{y}(E)=\{y\}\),
hence \(E\) is not a group if \(|E|>1\).}}

One may evidently suppose \(S'=S\); the translation \(t\) in question induces
on \(X_{0}\) a translation \(t_{0}\) of \(X_{0}\), which is therefore an
automorphism since \(X_{0}\) is a group. One concludes by flatness
(SGA~1 III~4.2).%
{\renewcommand{\thefootnote}{(67)}\nde{Since \(X\) and \(X_{0}\) have the same
underlying topological space and \(t_{0}\) is an automorphism, \(t\) is a
homeomorphism, hence an \emph{affine} morphism, \cf Exp.~VI\(_{\mathrm{B}}\),
2.9.1 or EGA~IV\(_{4}\), 18.12.7.1. It therefore suffices to see that if \(J\)
is a nilpotent ideal of a ring \(\Lambda\), and \(\phi\colon A\to B\) a morphism
of \(\Lambda\)-algebras, with \(B\) \emph{flat} over \(\Lambda\), such that
\(\phi\otimes_{\Lambda}(\Lambda/J)\) is bijective, then \(\phi\) is bijective.
By the ``nilpotent Nakayama lemma'', \(\phi\) is surjective; moreover, \(B\)
being flat over \(\Lambda\), one also has
\(\Ker(\phi)\otimes_{\Lambda}(\Lambda/J)=0\), whence \(\Ker(\phi)=0\), hence
\(\phi\) is bijective.}}

No longer assuming now that \(X\) has a section over \(S\), suppose that there
exists an \(S'\to S\) such that \(X_{S'}\) has a section over \(S'\). Then
\(X_{S'}\) is an \(S'\)-group by what one has just seen; let us consider its
unit section \(e'\). The inverse image of \(e'\) under
\(\mathrm{pr}_{i}\colon S'\times_{S}S'\to S'\) (\(i=1,2\)) is the unit section
of \(X_{S''}\) for the group law inverse image of \(P_{S'}\) under
\(\mathrm{pr}_{i}\). But as \(P\) is ``defined over \(S\)'', these two group
laws coincide, hence also their unit section. One therefore has
\(\mathrm{pr}_{1}^{*}(e')=\mathrm{pr}_{2}^{*}(e')\).

If \(S'\to S\) is a descent morphism (\cf Exp.~IV \No 2), there will exist a
section of \(X\) giving \(e'\) by extension of the base, and one will be done.
As \(X_{X}\) has a section over \(X\) (the diagonal section), one sees that it
now suffices to prove that \(X\to S\) is a descent morphism. Now it is flat and
surjective, and quasi-compact or locally of finite presentation, hence covering
for (\fpqc), hence a descent morphism (Exp.~IV, \No 6).
\end{proof}

\begin{remark*}
In fact the hypothesis that \(X\to S\) be quasi-compact or locally of finite
presentation is superfluous, by virtue of the following result, which the
reader will prove as an exercise on Exposé~IV:

Under the conditions of the text on \(S\) and \(S_{0}\), if \(X\to S\) is a
flat morphism and \(X_{0}\to S_{0}\) a morphism covering for (\fpqc), then
\(X\to S\) is a descent morphism.
\end{remark*}

\begin{lemma}{3.4}\label{III.3.4}
For two admissible composition laws on \(X\) to be equivalent (\cf 1.3.5), it
is necessary and sufficient that they be deduced from one another by an
automorphism of \(X\) over \(S\) inducing the identity on \(X_{\mathcal{J}}\).
\end{lemma}

\begin{proof}
\renewcommand{\qedsymbol}{}
Indeed, the morphisms constructed in~1.3.1 are exactly those of the preceding
statement (by~0.7).%
{\renewcommand{\thefootnote}{(68)}\nde{Indeed, by the proof of~0.7, the
\(S\)-endomorphisms of \(X\) inducing the identity on \(X_{\mathcal{J}}\) are
the automorphisms \(m\cdot\mathrm{id}_{X}\), for \(m\) running through
\(M(X)=\Hom_{S}(X,M)\) (for every \(S'\to S\) and \(x\in X(S')\), one has
\((m\cdot\mathrm{id}_{X})(x)=m(x)\cdot x\)). Now, by the proof of~3.1, each
\(m\colon X\to M\) factors uniquely as a morphism \(h\) from \(Y=X^{+}\) to
\(M\), and therefore \(m\cdot\mathrm{id}_{X}\) is the automorphism \(u_{h}\)
introduced in~1.3.1. The lemma then follows from the definition of the
equivalence, \cf 1.3.4 and~1.3.5.}}
\end{proof}

Taking all the preceding results into account, Proposition~1.3.6 gives:

\begin{theorem}{3.5}\label{III.3.5}
Let \(S\) be a scheme, \(\mathcal{I}\) and \(\mathcal{J}\) two ideals on \(S\)
such that \(\mathcal{I}\supset\mathcal{J}\),
\(\mathcal{I}\cdot\mathcal{J}=0\), and \(S_{0}\) and \(S_{\mathcal{J}}\) the
closed subschemes of \(S\) which they define. Let \(X\) be an \(S\)-scheme flat
over \(S\) (and locally of finite presentation or quasi-compact over \(S\)),
and \(X_{0}\) and \(X_{\mathcal{J}}\) the schemes obtained by base change.
Suppose \(X_{\mathcal{J}}\) endowed with a structure of
\(S_{\mathcal{J}}\)-group, and denote by \(L_{0}\) the \(S_{0}\)-functor in
abelian groups defined by the formula
\[
\Hom_{S_{0}}(T,L_{0})
  =\Hom_{\mathcal{O}_{T}}\bigl(
    \omega^{1}_{X_{0}/S_{0}}\otimes_{\mathcal{O}_{S_{0}}}\mathcal{O}_{T},\,
    \mathcal{J}\otimes_{\mathcal{O}_{S_{0}}}\mathcal{O}_{T}\bigr)
\]
on which \(X_{0}\) acts via its adjoint representation.
\begin{enumeratei}
\item For there to exist an \(S\)-group structure on \(X\) inducing the given
structure on \(X_{\mathcal{J}}\), it is necessary and sufficient that the
following conditions be satisfied:
\begin{enumerate}[label={($i_{\arabic*}$)}]
\item There exists a morphism of \(S\)-schemes \(X\times_{S}X\to X\) inducing
the group law of \(X_{\mathcal{J}}\).
\item A certain obstruction class belonging to \(H^{3}(X_{0},L_{0})\) (defined
canonically by the datum of \(X\) and of the group law of \(X_{\mathcal{J}}\))
is zero.
\end{enumerate}
\item If the conditions of~(i) are satisfied, the set \(E\) of group laws on
\(X\) inducing the given law of \(X_{\mathcal{J}}\) is a principal homogeneous
set under \(Z^{2}(X_{0},L_{0})\), and \(E\) modulo the \(S\)-automorphisms of
\(X\) inducing the identity on \(X_{\mathcal{J}}\), is a principal homogeneous
set under \(H^{2}(X_{0},L_{0})\).
\end{enumeratei}
\end{theorem}

\begin{proof}
\renewcommand{\qedsymbol}{}
{\renewcommand{\thefootnote}{(69)}\nde{One has added what follows.}}
Indeed, every morphism of \(S\)-schemes \(f\colon X\times_{S}X\to X\) inducing
the group law of \(X_{\mathcal{J}}\) is, by~3.2, an \emph{admissible}
composition law on \(X\); then, by~1.3.6 (i), the existence of an
\emph{associative} admissible composition law \(P\colon X\times_{S}X\to X\) is
equivalent to the vanishing of a certain class \(c(f)\in H^{3}(X_{0},L_{0})\),
and in this case, by~3.3, \(P\) is a group law. This proves~(i), and~(ii) then
follows from~3.3 and~1.3.6 (ii).
\end{proof}

\begin{remark}{3.5.1}\label{III.3.5.1}
{\renewcommand{\thefootnote}{(70)}\nde{One has added this remark, analogous
to~4.5.1, in order to introduce the notation \(\delta(\mu,\mu')\) (or
\(\delta(X,X')\)), used in~4.38; consequently, one has also added in~3.5 (ii)
above the part concerning \(E\) itself.}}
If \(\mu\), \(\mu'\) are group laws on \(X\) inducing the given law of
\(X_{\mathcal{J}}\), one therefore obtains a cocycle
\(\delta(\mu,\mu')\in Z^{2}(X_{0},L_{0})\), the sign convention chosen being
that \(\mu'=\delta(\mu,\mu')\cdot\mu\), \ie for every \(S'\to S\) and
\(x,y\in X(S')\),
\begin{equation*}
\mu'(x,y)=\delta(\mu,\mu')(x_{0},y_{0})\cdot\mu(x,y).
\end{equation*}%
{\renewcommand{\thefootnote}{(71)}\nde{One has conformed to the sign
conventions of the original, so as to have in~4.38 (5) the equality
\(\partial^{1}d(X,X')=\delta(X,X')\) (see also N.D.E.~(54)).}}
\end{remark}
